% Generated draft: requires mathematical and source review.
\documentclass[11pt]{article}
\usepackage[margin=25mm]{geometry}
\usepackage{fontspec}
\setmainfont{Latin Modern Roman}
\usepackage{amsmath,amssymb,mathtools,unicode-math}
\setmathfont{Latin Modern Math}
\usepackage{enumitem,needspace,xurl,hyperref}
\setlength{\parindent}{0pt}
\setlength{\parskip}{5pt}
\setlength{\emergencystretch}{3em}
\newcommand{\N}{\mathbb{N}}
\newcommand{\Z}{\mathbb{Z}}
\newcommand{\Q}{\mathbb{Q}}
\newcommand{\R}{\mathbb{R}}
\newcommand{\C}{\mathbb{C}}
\newcommand{\F}{\mathbb{F}}
\newcommand{\A}{\mathbb{A}}
\newcommand{\PP}{\mathbb{P}}
\newcommand{\E}{\mathbb{E}}
\newcommand{\eps}{\varepsilon}
\newcommand{\dd}{\,\mathrm{d}}
\begin{document}
\section*{UnsolvedMath ID: 30006329}
\textbf{Source problem code:} catalogue entry
\subsection{Definitions and mathematical statement}
The supplied fragment is headed ``Tempered Numerical Conjecture'' and labels the conjecture ``Conjecture 2, [2, 14.2.1].'' It chooses a coefficient field \(k\simeq\C\). Its explicit input is an everywhere-unramified tempered automorphic representation \(\pi\), together with a corresponding extended Langlands parameter, attributed to [6] and described by reference to [2, Chapter 2.6], of the displayed form
\[
\phi\colon\Gamma_F\longrightarrow\check G\times\mathbb G_m.
\]
Here \(\mathbb G_m\) denotes the multiplicative algebraic group. The fragment does not specify \(F\), define \(\Gamma_F\), identify the automorphic group or its dual \(\check G\), or give the conditions imposed on this extended parameter. These omissions prevent a self-contained definition of the class of representations under consideration.

The registry's supplied formulation additionally requires a hyperspherical dual pair. Neither that pair nor the meaning of hyperspherical duality is defined in the supplied excerpt. The proposed assertion is that, for every such pair and every everywhere-unramified tempered automorphic representation satisfying the intended hypotheses, the normalized automorphic period equals the normalized \(L\)-value prescribed by the corresponding extended parameter. Schematically, this is
\[
\underbrace{\text{normalized automorphic period of }\pi}_{\text{definition and normalization not supplied}}
=
\underbrace{\text{normalized }L\text{-value prescribed by }\phi}_{\text{definition and evaluation point not supplied}}.
\]
This display records the supplied intended target; it is not a recovered explicit formula. The excerpt neither defines the period nor specifies the \(L\)-function, its evaluation point, or the normalization of either side. Its disconnected final clause describes ``the latter terms'' as Artin \(L\)-functions of \(\Gamma_F\), normalized according to [2, Chapter 11], but the terms to which this refers are absent. Consequently, the complete equality and its precise quantifiers cannot be reconstructed from the supplied evidence. The normalization qualification must be retained even though the exposition invites readers to omit the ``norm'' decoration on a first reading.
\subsection{Short English statement}
The supplied conjectural target equates a normalized automorphic period with a normalized \(L\)-value for an everywhere-unramified tempered representation and a hyperspherical dual pair. The excerpt supplies only part of the input and a disconnected normalization reference, so the exact formula and hypotheses remain unresolved.
\subsection{Sources}
\begin{itemize}
\item Supplied local UnsolvedMath registry record, permanent ID \(30006329\), exact source problem code \texttt{catalogue entry}, title ``The Tempered Numerical Period Conjecture.'' The record assigns proposed year 2025 and supplies no individual proposer. Its status is recorded as open, and its statement status as \texttt{corrected\_verified}; these are supplied registry labels, not independent verification by this draft. The registry also reports that the official report was checked and the notation repaired; that verification claim is supplied provenance.
\item \emph{Arbeitsgemeinschaft: Relative Langlands Duality} (2025), Oberwolfach Reports, pp.~811--882, DOI \texttt{10.4171/owr/2025/18}. Supplied source and literature link: \url{https://doi.org/10.4171/owr/2025/18}. The extracted fragment carries page number 875 and the heading ``Tempered Numerical Conjecture.'' The supplied background identifies the report as documentation of workshop discussions and open problems.
\item \emph{Relative Langlands Duality}, supplied literature link: \url{https://arxiv.org/abs/2409.04677}. No authors or explicit identification with the excerpt's numbered bibliography entries are supplied in this record.
\item The fragment cites [2, 14.2.1] for Conjecture 2, [2, Chapter 2.6] for extended parameters, [6] for provision of the parameter, and [2, Chapter 11] for normalization of Artin \(L\)-functions. The bibliography identifying [2] and [6] is not included in the supplied evidence.
\item Literature assessment dated exactly \texttt{2026-08-22}: the supplied curation labels the problem open, describes the fragment as stating the input to the Tempered Numerical Conjecture, and says that the report contains examples and extensions while the general period-equals-\(L\)-value formula remains unresolved. This is supplied dated curation, not independently verified. The background expressly requests independent verification of time-sensitive status and cited claims.
\end{itemize}
\end{document}
